\section*{Разбор задачи из тестового варианта №4}
\addcontentsline{toc}{section}{Разбор задачи из тестового варианта №4}
\begin{Task}{Ожидания функций от независимых величин и смешанные распределения}{}
Случайные величины $X$ и $Y$ независимы, $X$ имеет функцию распределения
\[
F(x) =
\begin{cases}
0, & x < 0; \\
\frac{x^2+2x+1}{15}, & x \in [0, 2); \\
\frac{3}{4}, & x \in [2, 3); \\
1, & x \geqslant 3,
\end{cases}
\]
а $Y$ имеет равномерное распределение на $(1, 2)$. Найдите $\E Y^X$.
\end{Task}

\begin{Solution}
Анализ структуры задачи указывает на необходимость вычисления математического ожидания нелинейной функции от двух независимых случайных величин, одна из которых ($X$) имеет смешанное распределение. Прямое вычисление двойного интеграла Лебега-Стилтьеса неоптимально. Фундаментальный подход заключается в последовательном интегрировании (использовании формулы полного математического ожидания) с предварительной декомпозицией распределения $X$.

\begin{MethodBox}[Оптимизация: Внутреннее ожидание]
В силу независимости $X$ и $Y$, вычисление $\E[Y^X]$ сводится к нахождению математического ожидания вспомогательной функции $h(x)$, определяемой как условное (в данном случае — параметрическое) ожидание по $Y$ при фиксированном $X=x$:
\[
h(x) = \E[Y^x] = \int_{-\infty}^{+\infty} y^x p_Y(y) \, dy.
\]
Искомая величина вычисляется как интеграл по мере $X$: $\E[Y^X] = \E[h(X)]$.
\end{MethodBox}

\textbf{Шаг 1. Вычисление функции $h(x)$.} \\
Случайная величина $Y \sim U(1, 2)$ абсолютно непрерывна и имеет плотность $p_Y(y) = 1$ при $y \in (1, 2)$ и $p_Y(y) = 0$ иначе. Вычислим внутреннее ожидание:
\[
h(x) = \int_{1}^{2} y^x \cdot 1 \, dy = \left. \frac{y^{x+1}}{x+1} \right|_{1}^{2} = \frac{2^{x+1} - 1}{x+1}.
\]
\textit{Замечание:} Из условия следует, что $X \geqslant 0$ почти наверное, следовательно, знаменатель $x+1 \neq 0$ и функция $h(x)$ определена корректно.

\textbf{Шаг 2. Декомпозиция распределения $X$.} \\
Функция распределения $F(x)$ имеет как интервалы непрерывного роста, так и точки разрыва первого рода (скачки). Представим распределение $X$ в виде суммы дискретной и абсолютно непрерывной компонент.

\textit{Анализ дискретной части (точки разрыва):}
Величина скачка в точке $x_k$ определяет вероятность $\P(X = x_k) = F(x_k) - F(x_k - 0)$.
Исследуем границы интервалов:
\begin{itemize}
    \item В точке $x = 0$: $F(0) = \frac{1}{15}$, $F(0-0) = 0$. Скачок: $\P(X = 0) = \frac{1}{15}$.
    \item В точке $x = 2$: $F(2) = \frac{3}{4} = \frac{15}{20}$, $F(2-0) = \frac{2^2+2(2)+1}{15} = \frac{9}{15} = \frac{12}{20}$. Скачок: $\P(X = 2) = \frac{15}{20} - \frac{12}{20} = \frac{3}{20}$.
    \item В точке $x = 3$: $F(3) = 1$, $F(3-0) = \frac{3}{4}$. Скачок: $\P(X = 3) = 1 - \frac{3}{4} = \frac{1}{4}$.
\end{itemize}

\textit{Анализ абсолютно непрерывной части:}
На интервалах дифференцируемости $F(x)$ найдем плотность $p_X(x) = F'(x)$:
\begin{itemize}
    \item При $x \in (0, 2)$: $p_X(x) = \left( \frac{x^2+2x+1}{15} \right)' = \frac{2(x+1)}{15}$.
    \item На остальных интервалах $(-\infty, 0) \cup (2, 3) \cup (3, +\infty)$ производная $F'(x) = 0$.
\end{itemize}

\begin{WarningBox}[Проверка условия нормировки]
Для смешанного распределения сумма всех вероятностей скачков и интеграла от непрерывной плотности должна строго равняться единице.
\vspace{-0.5em}
\begin{align*}
    \sum_{k} \P(X=x_k) + \int_{\mathbb{R}} p_X(x) \, dx &= \left( \frac{1}{15} + \frac{3}{20} + \frac{1}{4} \right) + \int_{0}^{2} \frac{2(x+1)}{15} \, dx \\
    &= \frac{4 + 9 + 15}{60} + \left. \frac{(x+1)^2}{15} \right|_{0}^{2} \\
    &= \frac{28}{60} + \left( \frac{9}{15} - \frac{1}{15} \right) = \frac{7}{15} + \frac{8}{15} = 1.
\end{align*}
\vspace{-1em}

Декомпозиция проведена корректно, потери массы вероятности нет.
\end{WarningBox}

\textbf{Шаг 3. Синтез (расчет итогового ожидания).} \\
Применим формулу интегрирования по мере Стилтьеса:
\[
\E[h(X)] = \sum_{k} h(x_k)\P(X=x_k) + \int_{0}^{2} h(x)p_X(x) \, dx.
\]

Вычислим вклад дискретной компоненты:
\begin{align*}
    h(0)\P(X=0) &= \frac{2^1 - 1}{0+1} \cdot \frac{1}{15} = 1 \cdot \frac{1}{15} = \frac{1}{15}, \\
    h(2)\P(X=2) &= \frac{2^3 - 1}{2+1} \cdot \frac{3}{20} = \frac{7}{3} \cdot \frac{3}{20} = \frac{7}{20}, \\
    h(3)\P(X=3) &= \frac{2^4 - 1}{3+1} \cdot \frac{1}{4} = \frac{15}{4} \cdot \frac{1}{4} = \frac{15}{16}.
\end{align*}
Сумма дискретных слагаемых:
\[
S_d = \frac{1}{15} + \frac{7}{20} + \frac{15}{16} = \frac{16 + 84 + 225}{240} = \frac{325}{240} = \frac{65}{48}.
\]

Вычислим вклад абсолютно непрерывной компоненты. Обратите внимание на алгебраическое сокращение $(x+1)$, которое существенно упрощает интеграл:
\begin{align*}
    I_c &= \int_{0}^{2} \left( \frac{2^{x+1} - 1}{x+1} \right) \cdot \frac{2(x+1)}{15} \, dx = \frac{2}{15} \int_{0}^{2} (2^{x+1} - 1) \, dx \\
    &= \frac{2}{15} \left. \left( \frac{2^{x+1}}{\ln 2} - x \right) \right|_{0}^{2} = \frac{2}{15} \left[ \left( \frac{8}{\ln 2} - 2 \right) - \left( \frac{2}{\ln 2} - 0 \right) \right] \\
    &= \frac{2}{15} \left( \frac{6}{\ln 2} - 2 \right) = \frac{4}{5 \ln 2} - \frac{4}{15}.
\end{align*}

Окончательное суммирование:
\begin{align*}
    \E[Y^X] &= S_d + I_c = \frac{65}{48} + \left( \frac{4}{5 \ln 2} - \frac{4}{15} \right) \\
    &= \left( \frac{65 \cdot 5 - 4 \cdot 16}{240} \right) + \frac{4}{5 \ln 2} \\
    &= \frac{325 - 64}{240} + \frac{4}{5 \ln 2} = \frac{261}{240} + \frac{4}{5 \ln 2} = \frac{87}{80} + \frac{4}{5 \ln 2}.
\end{align*}
\end{Solution}

\begin{Answer}
$\E Y^X = \frac{87}{80} + \frac{4}{5 \ln 2}$
\end{Answer}
